\chapter{A Concrete Naming Certificate for $\RealIntervalCauchyFilterUp$}
\label{ch:concrete-instances-real-interval-cauchy-filter-namecert}
\origin{human}

Following Bishop and Bridges constructive analysis and Bourbaki's finite-neighbourhood reading of filters, the $\RealIntervalCauchyFilterUp$ packet records an interval-indexed Cauchy filter around displayed real names as finite BEDC NameCert data. It sits between the terminal real-name seal of \autoref{ch:concrete-instances-real-namecert}, the filter surface of \autoref{ch:concrete-instances-cauchyfilter-namecert}, the dyadic tolerance core of \autoref{ch:concrete-instances-dyadicratcore-namecert}, the finite stream schedule of \autoref{ch:concrete-instances-streamname-namecert}, the regular rational readback of \autoref{ch:concrete-instances-regseqrat-namecert}, and the metric comparison row of \autoref{ch:concrete-instances-metric-namecert}. The packet names only the finite interval-filter route by which Real and Completion consumers inspect a Cauchy neighbourhood; it does not select a limit, quotient filter, host real equality, open-cover compactness principle, countable-choice selector, or ambient $\RealUp$ completeness witness.

\begin{definition}[RealIntervalCauchyFilter carrier]
\label{def:real-interval-cauchy-filter-carrier}
A $\RealIntervalCauchyFilterUp$ carrier is a finite $\BHist$ packet
$$
I=(R,F,D,W,Q,M,H,C,P,N)
$$
with the following displayed rows:
\begin{enumerate}
\item $R$ is the $\RealUp$ source row naming the real endpoint or center data whose interval neighbourhoods are inspected.
\item $F$ is the $\CauchyFilterUp$ finite filter-basis row, listing the accepted interval-neighbourhood requests and their tail order.
\item $D$ is the $\DyadicRatCoreUp$ tolerance row recording radii, refinements, and dyadic containment comparisons for the inspected intervals.
\item $W$ is the $\StreamNameUp$ finite-window row that opens only the stream prefixes needed by the filter request.
\item $Q$ is the $\RegSeqRatUp$ readback row for the displayed windows and dyadic tolerances.
\item $M$ is the $\MetricUp$ comparison row that records interval containment and diameter checks before any Real-facing handoff is exposed.
\item $H$ records componentwise $\hsame$ transport for the Real source, filter basis, dyadic tolerance row, stream window, regular readback, metric comparison, replay, provenance, and local naming.
\item $C$ records displayed $\Cont$ replay from an interval-filter request through $R,F,D,W,Q,M,H$.
\item $P$ is the $\Pkg$ provenance over $\RealIntervalCauchyFilterUp$, $\RealUp$, $\CauchyFilterUp$, $\DyadicRatCoreUp$, $\StreamNameUp$, $\RegSeqRatUp$, $\MetricUp$, $\BHist$, $\hsame$, $\Cont$, and $\NameCert$ rows.
\item $N$ is the local $\NameCert_{\RealIntervalCauchyFilterUp}$ row.
\end{enumerate}
The carrier has no coordinate for quotient filters, selected limits, countable choice, host real equality, open-cover compactness, or ambient $\RealUp$ completeness.
\end{definition}

\begin{theorem}[RealIntervalCauchyFilter NameCert obligations]
\label{thm:real-interval-cauchy-filter-namecert-obligations}
For every accepted $\RealIntervalCauchyFilterUp$ carrier $I=(R,F,D,W,Q,M,H,C,P,N)$, the local $\NameCert_{\RealIntervalCauchyFilterUp}$ surface consists exactly of Real source admission through $R$, finite Cauchy filter-basis exposure through $F$, dyadic interval tolerance exposure through $D$, stream-window exposure through $W$, regular rational readback through $Q$, metric interval comparison through $M$, componentwise $\hsame$ transport through $H$, displayed $\Cont$ replay through $C$, $\Pkg$ provenance through $P$, and local naming through $N$. A consumer therefore reads interval Cauchy-filter data only through the finite route
$$
R,F,D,W,Q,M \longmapsto H,C,P,N .
$$
\end{theorem}
\begin{proof}
Project the accepted carrier of \autoref{def:real-interval-cauchy-filter-carrier}. The rows $R$ and $F$ expose the Real source and finite filter-basis request before any interval read is available.

The dyadic tolerance row $D$, stream window $W$, regular readback $Q$, and metric comparison row $M$ then display the interval containment route. These rows make the inspected neighbourhood finite and visible rather than a hidden quotient filter or selected completion object.

Finally the structural rows $H,C,P,N$ transport, replay, source, and name the same carrier. Since the carrier has no coordinate for quotient filters, selected limits, countable choice, host real equality, open-cover compactness, or ambient Real completeness, the listed rows exhaust the NameCert obligations.
\end{proof}

\begin{theorem}[RealIntervalCauchyFilter tail basis]
\label{thm:real-interval-cauchy-filter-tail-basis}
Every accepted $\RealIntervalCauchyFilterUp$ tail read factors through the same finite filter-basis row $F$, dyadic tolerance row $D$, $\StreamNameUp$ window $W$, $\RegSeqRatUp$ readback $Q$, and $\MetricUp$ comparison row $M$ before $H,C,P,N$ may replay or name the read. In particular, interval refinement is a displayed basis operation inside $I=(R,F,D,W,Q,M,H,C,P,N)$ and not an appeal to an ambient neighbourhood filter or countable family of representatives.
\end{theorem}
\begin{proof}
Start with the NameCert obligations of \autoref{thm:real-interval-cauchy-filter-namecert-obligations}. They force every accepted tail read to expose the filter-basis row $F$ and the dyadic interval row $D$ before the route can reach stream or rational readback data.

The finite window $W$ and regular readback $Q$ then inspect only the prefixes and rational values listed by the carrier. The metric comparison row $M$ records containment and diameter checks on that same displayed basis.

The rows $H,C,P,N$ preserve the basis route under transport, continuation replay, provenance, and local naming. Thus no ambient neighbourhood filter, quotient representative, selected limit, or countable-choice schedule is exported by the tail read.
\end{proof}

\begin{theorem}[RealIntervalCauchyFilter Real handoff]
\label{thm:real-interval-cauchy-filter-real-handoff}
Every Real-facing consumer of an accepted $\RealIntervalCauchyFilterUp$ packet must consume the displayed route
$$
\RealUp,\ \CauchyFilterUp,\ \DyadicRatCoreUp,\ \StreamNameUp,\ \RegSeqRatUp,\ \MetricUp
$$
through the carrier rows $R,F,D,W,Q,M,H,C,P,N$. The handoff supplies a finite interval-filter witness for Real and Completion neighbours, but it exports no quotient filter, selected limit, host real equality, open-cover compactness theorem, countable-choice selector, or ambient $\RealUp$ completeness principle.
\end{theorem}
\begin{proof}
Apply the tail-basis theorem, \autoref{thm:real-interval-cauchy-filter-tail-basis}. It keeps every interval-filter read on the displayed basis row $F$, dyadic tolerance row $D$, stream window $W$, regular readback $Q$, and metric comparison row $M$.

The Real source row $R$ is therefore consumed together with the same finite basis route. No consumer can detach the terminal Real read from the filter, dyadic, stream, rational, and metric rows that make the interval request visible.

Finally use the structural rows $H,C,P,N$ from the NameCert obligations. They replay and name the same finite packet, so the handoff cannot introduce quotient filters, selected limits, host real equality, open-cover compactness, countable choice, or ambient Real completeness.
\end{proof}

\falsifiablePrediction{Every accepted $\RealIntervalCauchyFilterUp$ read exposes a RealUp source row, a finite CauchyFilterUp basis row, a DyadicRatCoreUp interval tolerance row, a StreamNameUp finite window, a RegSeqRatUp readback row, a MetricUp comparison row, componentwise hsame transport, displayed Cont replay, Pkg provenance, and local naming before any Real-facing handoff is visible.}

\independenceWitness{real, cauchyfilter, dyadicratcore, streamname, regseqrat, metric: $\RealIntervalCauchyFilterUp$ is not bijective to any listed sibling because it jointly binds Real interval source data, finite filter-basis rows, dyadic tolerance rows, finite stream windows, regular rational readback, metric interval comparisons, transport, replay, provenance, and local naming.}

\closureat{RealIntervalCauchyFilterUp}{seedStr}

\begin{closurestatus}{\RealIntervalCauchyFilterUp}
  \constructivestory{A $\RealIntervalCauchyFilterUp$ packet is generated from a RealUp source row, a finite CauchyFilterUp basis row, a DyadicRatCoreUp interval tolerance row, a StreamNameUp finite window, a RegSeqRatUp readback row, a MetricUp comparison row, componentwise hsame transport, displayed Cont replay, Pkg provenance, and a local NameCert row.}
  \theoryclosure{\seedClosure}
  \scopeclosed{\autoref{def:real-interval-cauchy-filter-carrier} and \autoref{thm:real-interval-cauchy-filter-namecert-obligations} seed the finite interval Cauchy-filter carrier over RealUp, CauchyFilterUp, DyadicRatCoreUp, StreamNameUp, RegSeqRatUp, MetricUp, BHist, hsame, Cont, Pkg, and NameCert rows.}
  \formalstatus{\unformalizedV}
  \bridgestatus{none}
  \notclaimed{This chapter does not claim quotient filters, selected limits, countable choice, host real equality, open-cover compactness, ambient Real completeness, Lean verification, or bridge checking.}
  \upgradepath{Obligation closure requires checked carrier admission, finite filter-basis exactness, dyadic interval tolerance stability, StreamName window admission, RegSeqRat readback, Metric comparison nonescape, replay stability, and local naming for BEDC.Derived.RealIntervalCauchyFilterUp.}
\end{closurestatus}
